\chapter{Bevezetés}
\label{ch:bevezetes}
% Terjedelem: kb. 3–4 oldal (~5 000 karakter)
% 15. pont: a témaválasztás indokai + milyen elméleti és gyakorlati kérdéseket vizsgál a dolgozat.
% 17. pont: összefüggő szöveg, NE felsorolás!

\section{A témaválasztás indoklása}
% - Az online / digitális vizsgáztatás elterjedése (COVID-19 után különösen), a felsőoktatás igényei.
% - Informatikai képzésben a programozási feladatok kézi javítása időigényes és nem egységes.
% - Meglévő rendszerek (Moodle Quiz, CodeRunner stb.) hiányosságai: kódfuttatás biztonsága,
%   valós idejű felügyelet hiánya, összetett beállítás.
% - Saját motiváció (saját vélemény is kell – 19. pont).
\todo{Szöveg}

\section{Célkitűzések}
% Egy olyan webes prototípus elkészítése, amely:
%   vegyes feladattípusokat kezel (kód, egy/több választós, igaz–hamis, rövid válasz, esszé),
%   a programozási feladatokat biztonságos Docker-sandboxban, tesztesetekkel automatikusan pontozza,
%   valós időben jelzi az oktatónak, ha a hallgató elhagyja a vizsgaablakot,
%   oktatói felületet ad a kézi javításhoz és az eredmények exportálásához.
% Vizsgált kérdések: hogyan lehet nem megbízható kódot biztonságosan futtatni?
% Milyen pontozási szabályok igazságosak? Mennyit ér a böngészőalapú felügyelet?
\todo{Szöveg}

\section{A dolgozat felépítése}
% Röviden, fejezetenként (\ref{ch:szakirodalom}. fejezet ..., \ref{ch:technologiak}. fejezet ...).
\todo{Szöveg}
